\par\medskip
\begin{center}\hypertarget{native-164-article117}{}\pdfbookmark[2]{117.}{article-117}117.\end{center}
Positio terrae hoc modo e tabulis computata ad terrae centrum referenda est, locus observatus autem corporis coelestis ad punctum in terrae superficie spectat: huic dissensui tribus modis remedium afferre licet. Potest scilicet vel observatio ad centrum terrae reduci, sive a parallaxi liberari; vel locus heliocentricus terrae ad locum ipsum observationis reduci, quod efficitur, si loco Solis e tabulis computato parallaxis rite applicatur; vel denique utraque positio ad punctum aliquod tertium transferri, quod commodissime in intersectione radii visus cum plano eclipticae assumitur: observatio ipsa tunc immutata manet, reductionemque loci terrae ad hoc punctum in art. 72 docuimus. Methodus prima adhiberi nequit, nisi corporis coelestis distantia a terra proxime saltem nota fuerit: tunc autem satis commoda est, praesertim quoties observatio in ipso meridiano instituta est, ubi sola declinatio parallaxi afficitur. Ceterum praestabit, hanc reductionem loco observato immediate applicare, antequam transformationes art. praec. adeantur. Si vero distantia a terra penitus adhuc incognita est, ad methodum secundam vel tertiam confugiendum est, et quidem illa in usum vocabitur, quoties aequator pro plano fundamentali accipitur, tertia autem praeferetur, quoties omnes positiones ad eclipticam referre placuit.
